\section{Authentic teaching and religious submission}

\emph{Lumen gentium} 25 requires a religious submission of intellect and will
to authentic papal teaching even when the pope does not speak ex cathedra, and
it directs attention to the document's character, repetition, and manner of
speaking. This response is neither divine faith nor mere external silence. It
is a real presumption of truth and a loyal willingness to be taught, arising
from Christ's ordering of the Church.

Non-definitive teaching can contain contingent elements, incomplete
formulations, or prudential applications capable of development and correction.
Calling it reformable does not mean that a private reader may suspend it at
will. \emph{Donum veritatis} distinguishes a serious difficulty, pursued with
competence, patience, and ecclesial loyalty, from organized pressure, appeal to
public opinion, or a parallel teaching office. Difficulty should be stated
exactly: the proposition, authority, evidence, prior teaching, and feared
consequence. A Catholic continues to avoid contempt, scandal, and the claim
that personal judgment is self-authenticating.

The opposite fault is to demand internal certainty by force. Authority can
oblige a posture and conduct; it cannot make a person perceive an argument that
he does not yet perceive. Pastors and theologians serve communion by clarifying
the object, removing caricatures, distinguishing levels, and permitting honest
questions within their proper bounds.

